\section{Computations}\label{tc:sec:comp}

The programs named below accompany \cite{PaperB}, in its directory \texttt{code/}; the file \texttt{README.md} there states, for each program, what it computes and its expected output. Proposition~\ref{pm:prop:cert} is the only statement of this note that is computer-assisted, and its consequences for the holonomy are also proved by hand (Proposition~\ref{pm:prop:hol}); the other computations confirm results proved by hand.

\emph{The lift certificate} (Proposition~\ref{pm:prop:cert}). The program \texttt{lift\_cert.py} computes the node sets $\mathcal R_k$ as specified in Section~\ref{pm:sec:cert}, with the parameters $(f,n_{\max},T_{\max})$ of the proposition, up to the depth at which the repetition stated there occurs for (1) and (2), and up to depth $105$ for (3). Each new set is compared with all the sets before it; for (3) the comparison is through SHA-256 digests of the sorted lists of nodes, and the $106$ digests are pairwise different. The consequences of Proposition~\ref{pm:prop:cert} for the holonomy are also proved by hand (Proposition~\ref{pm:prop:hol}), and no other proof depends on this computation.

\emph{The two-copy window graph} (Theorem~\ref{r30:thm:window} and Proposition~\ref{pm:prop:twocopy}). For $f=30$ and all $(w,M)$ with $w+M\le6$, and for $f=45,75,135$ and $(w,M)\in\{(1,0),(2,0),(3,0),(4,0),(1,1),(2,1),(3,1),(1,2)\}$, the graph $\Gamma_{w,M}$ built from the cell rule has an eventual image of $2^{3w+2}4^M$ states, which is strongly connected and is reached by every walk of length $\lceil w/2\rceil$.

\emph{The holonomy behind the front} (Lemma~\ref{pm:lem:theta} and Proposition~\ref{pm:prop:odd}). For $135$ the graphs $\Gamma_{w,M}$ are built directly from the cell rule for $2\le w\le7$ with $M=0$, and for $(w,M)=(2,1)$, $(3,1)$, $(4,1)$, $(5,1)$, $(6,1)$, $(2,2)$, $(3,2)$, $(4,2)$, $(2,3)$. In every case the eventual image has $2^{3w+2}4^M$ states, the graph $G_w$ at the step $y[w]$, obtained from it by forgetting the layer $y[w]$, is strongly connected, and the values $\Psi_w(\varsigma)$ of its closed walks are all $32$ elements of $(\Z/2)^5$. The same holds for the other seven non-affine left-permutive rules for $2\le w\le4$, $M=0$. For the affine rules $90$, $105$ and $150$ with $2\le w\le4$, $M=0$, the values of all closed walks generate a group of order at most $8$ (order $8$ for $90$, whose graph is not strongly connected; order $2$ or $4$ for $105$ and $150$). The supplementary material contains four programs for this, which use the cell rule only and neither the layer equation nor the conditions $(C_j)$. The program \texttt{ref\_check.py} enumerates all window states, iterates to the eventual image, and finds the exact set of values $\Psi_w$ of the closed walks at a vertex by a search in $G_w\times(\Z/2)^5$; with the option \texttt{flip} it tests that the top flips and the swap of Lemma~\ref{pm:lem:theta} are automorphisms of $G_w$ that change $\psi_w$ as stated there. The program \texttt{ref\_fast.py} builds the eventual image by a breadth-first search from one of its states, counts the strongly connected components, and computes the group generated by the values $\Psi_w$ from the potentials of a spanning tree; it is used for the larger cases $w=6,7$ with $M=0$ and $(w,M)=(5,1)$, $(6,1)$, $(4,2)$. The program \texttt{ref\_affine.py} does the same over all components for the affine rules, and \texttt{ref\_base.py} checks that the history below $2$ of Section~\ref{pm:sec:behind} for $135$ is a closed walk of length $3$ in $G_2$ with $\Psi_2=(1,0,1,1,1)$ for $M\le3$, and that sampled closed walks of $G_{d-1}$ with $\ell=(0,0)$ lift to closed walks of $G_d$ of the same length.
